\subsection{REDUCED DECOMPOSITIONS IN A COXETER GROUP}

Suppose that \((W,S)\) is a Coxeter system. Let T be the set of conjugates in W of elements of S. For any finite sequence \(\mathbf{s}=(s_{1},\ldots,s_{q})\) of elements of S, denote by \(\Phi(\mathbf{s})\) the sequence \((t_{1},\ldots,t_{q})\) of elements of T defined by

\[
t _ {j} = \left(s _ {1} \dots s _ {j - 1}\right) s _ {j} \left(s _ {1} \dots s _ {j - 1}\right) ^ {- 1} \quad \text { for } 1 \leqslant j \leqslant q.\tag{11}
\]

Then \(t_1 = s_1\) and \(s_1 \ldots s_q = t_q t_{q-1} \ldots t_1\) . For any element \(t \in \mathbf{T}\) , denote by \(n(\mathbf{s}, t)\) the number of integers \(j\) such that \(1 \leqslant j \leqslant q\) and \(t_j = t\) . Finally, put

\[
R = \{1, - 1 \} \times T.
\]

Lemma 1. (i) Let \(w \in \mathbf{W}\) and \(t \in \mathbf{T}\) . The number \((-1)^{n(\mathbf{s},t)}\) has the same value \(\eta(w,t)\) for all sequences \(\mathbf{s} = (s_1, \ldots, s_q)\) of elements of \(S\) such that \(w = s_1 \ldots s_q\) .

\begin{enumerate}
\def\labelenumi{(\roman{enumi})}
\setcounter{enumi}{1}
\tightlist
\item
  For \(w \in \mathbf{W}\) , let \(\mathrm{U}_w\) be the map from \(\mathbf{R}\) to itself defined by
\end{enumerate}

\[
\mathrm{U} _ {w} (\varepsilon , t) = (\varepsilon . \eta (w ^ {- 1}, t), w t w ^ {- 1}) \quad (\varepsilon = \pm 1, t \in \mathrm{T}).\tag{12}
\]

The map \(w \mapsto \mathrm{U}_w\) is a homomorphism from \(\mathbf{W}\) to the group of permutations of \(\mathbb{R}\) .

For \(s \in S\) , define a map \(U_s\) from \(R\) to itself by

\[
\mathrm{U} _ {s} (\varepsilon , t) = \left(\varepsilon . (- 1) ^ {\delta_ {s, t}}, s t s ^ {- 1}\right) \quad (\varepsilon = \pm 1, t \in \mathrm{T}),\tag{13}
\]

where \(\delta_{s,t}\) is the Kronecker symbol. It is immediate that \(U_{s}^{2} = Id_{R}\) , which shows that \(U_{s}\) is a permutation of R.

Let \(\mathbf{s} = (s_1, \ldots, s_q)\) be a sequence of elements of \(S\) . Put \(w = s_q \ldots s_1\) and \(U_{\mathbf{s}} = U_{s_q} \ldots U_{s_1}\) . We shall show, by induction on \(q\) , that

\[
\mathrm{U} _ {\mathbf {s}} (\varepsilon , t) = (\varepsilon . (- 1) ^ {n (\mathbf {s}, t)}, w t w ^ {- 1}).\tag{14}
\]

This is clear for \(q = 0,1\) . If \(q > 1\) , put \(\mathbf{s}' = (s_1,\dots ,s_{q - 1})\) and

\[
w ^ {\prime} = s _ {q - 1} \dots s _ {1}.
\]

Using the induction hypothesis, we obtain

\[
\begin{array}{r l} \mathrm{U} _ {\mathbf {s}} (\varepsilon , t) & = \mathrm{U} _ {s _ {q}} (\varepsilon . (- 1) ^ {n (\mathbf {s} ^ {\prime}, t)}, w ^ {\prime} t w ^ {\prime - 1}) \\ & = (\varepsilon . (- 1) ^ {n (\mathbf {s} ^ {\prime}, t) + \delta_ {s _ {q}, w ^ {\prime} t w ^ {\prime - 1}}}, w t w ^ {- 1}). \end{array}
\]

But \(\Phi(\mathbf{s}) = (\Phi(\mathbf{s}'), w'^{-1}s_q w')\) and \(n(\mathbf{s}, t) = n(\mathbf{s}', t) + \delta_{w'^{-1}s_q w', t}\) , proving formula (14).

Now let \(s, s' \in S\) be such that \(p = ss'\) has finite order \(m\) . Let \(\mathbf{s} = (s_1, \ldots, s_{2m})\) be the sequence of elements of \(S\) defined by \(s_j = s\) for \(j\) odd and \(s_j = s'\) for \(j\) even. Then \(s_{2m} \ldots s_1 = p^{-m} = 1\) and formula (11) implies that

\[
t _ {j} = p ^ {j - 1} s \quad \text { for } 1 \leqslant j \leqslant 2 m.\tag{15}
\]

Since \(p\) is of order \(m\) , the elements \(t_1, \ldots, t_m\) are distinct and \(t_{j+m} = t_j\) for \(1 \leqslant j \leqslant m\) . For all \(t \in \mathbf{T}\) , the integer \(n(\mathbf{s}, t)\) is thus equal to 0 or 2 and (14) shows that \(\mathrm{U_s} = \mathrm{Id_R}\) . In other words, \((\mathrm{U_sU_{s'}})^m = \mathrm{Id_R}\) . Thus, by the definition of Coxeter systems, there exists a homomorphism \(w \mapsto \mathrm{U}_w\) from \(W\) to the group of permutations of \(R\) such that \(U_s\) is given by the right-hand side of (13). Then \(U_w = U_s\) for every sequence \(s = (s_1, \ldots, s_q)\) such that \(w = s_q \ldots s_1\) and Lemma 1 follows immediately from (14).

Lemma 2. Let \(\mathbf{s} = (s_1, \ldots, s_q)\) , \(\Phi(\mathbf{s}) = (t_1, \ldots, t_q)\) and \(w = s_1 \ldots s_q\) . Let \(\mathrm{T}_w\) be the set of elements of \(\mathrm{T}\) such that \(\eta(w, t) = -1\) . Then \(\mathbf{s}\) is a reduced decomposition of \(w\) if and only if the \(t_i\) are distinct, and in that case \(\mathrm{T}_w = \{t_1, \ldots, t_q\}\) and \(\operatorname{Card}(\mathrm{T}_w) = l(w)\) .

Clearly \(T_{w} \subset \{t_{1}, \ldots, t_{q}\}\) . Taking s to be reduced, it follows that \(\operatorname{Card}(T_{w}) \leqslant l(w)\) . Moreover, if the \(t_{i}\) are distinct, then \(n(\mathbf{s}, t)\) is equal to 1 or 0 according to whether t does or does not belong to \(\{t_{1}, \ldots, t_{q}\}\) . It follows that \(T_{w} = \{t_{1}, \ldots, t_{q}\}\) and that \(q = \operatorname{Card}(T_{w}) \leqslant l(w)\) , which implies that s is reduced. Suppose finally that \(t_{i} = t_{j}\) with i \textless{} j. This gives \(s_{i} = us_{j}u^{-1}\) , with \(u = s_{i+1} \ldots s_{j-1}\) , hence

\[
w = s _ {1} \dots s _ {i - 1} s _ {i + 1} \dots s _ {j - 1} s _ {j + 1} \dots s _ {q}.
\]

This shows that s is not a reduced decomposition of w.

Lemma 3. Let \(w \in \mathbb{W}\) and \(s \in \mathbb{S}\) be such that \(l(sw) \leqslant l(w)\) . For any sequence \(\mathbf{s} = (s_1, \ldots, s_q)\) of elements of \(\mathbb{S}\) with \(w = s_1 \ldots s_q\) , there exists an integer \(j\) such that \(1 \leqslant j \leqslant q\) and

\[
s s _ {1} \dots s _ {j - 1} = s _ {1} \dots s _ {j - 1} s _ {j}.
\]

Let \(p\) be the length of \(w\) and \(w' = sw\) . By the Remark of no. 3, \(l(w') \equiv l(w) + 1 \mod 2\) . The hypothesis \(l(w') \leqslant l(w)\) and the relation

\[
| l (w) - l (w ^ {\prime}) | \leqslant l (w w ^ {\prime - 1}) = l (s) = 1
\]

thus leads to \(l(w') = p - 1\) . Choose a reduced decomposition

\[
\left(s _ {1} ^ {\prime}, \ldots , s _ {p - 1} ^ {\prime}\right)
\]

of \(w'\) and put \(\mathbf{s} = (s, s_1', \ldots, s_{p-1}')\) and \(\Phi(\mathbf{s}') = (t_1', \ldots, t_p')\) . It is clear that \(\mathbf{s}'\) is a reduced decomposition of \(w\) and that \(t_1' = s\) . The elements \(t_1', \ldots, t_p'\) being distinct by Lemma 2, we have \(n(\mathbf{s}', s) = 1\) . Since \(w\) is the product of the sequence \(\mathbf{s}\) , we have \(n(\mathbf{s}, s) \equiv n(\mathbf{s}', s) \mod 2\) by Lemma 1, hence \(n(\mathbf{s}, s) \neq 0\) . Consequently, \(s\) is equal to one of the elements \(t_j\) of the sequence \(\Phi(\mathbf{s})\) , hence the lemma.

Remark. The set \(T_{w}\) defined in Lemma 2 consists of the elements of the form \(w''sw''^{-1}\) corresponding to the triples \((w', w'', s) \in W \times W \times S\) such that \(w = w''sw'\) and \(l(w') + l(w'') + 1 = l(w)\) .

